% ==========================================================================
% Figure blocks for
%   Shakespeare with Ladders: Composition, Derived Power, and the Second Refusal of a Receiver Language
%
% Each block below is self-contained. Paste one straight into
% the manuscript at the point its panel belongs; no \input and
% no \captionPanel* macro is required.
%
% Preamble requirements:
%   \usepackage{graphicx}
%   \usepackage{amsmath,amssymb}
%   \graphicspath{{figures/}}      % if the paths below are bare
%   \usepackage{cleveref}          % for the \cref calls below
%
% These captions cross-reference the manuscript. The following
% labels must exist, or LaTeX will report an undefined control
% sequence (missing cleveref) or an undefined reference:
%   eq:local-sepcost
%   rem:res-inverted
% Replace the \cref calls with plain text if you are pasting
% these panels into a document that lacks those labels.
% ==========================================================================


% --------------------------------------------------------------------------
\begin{figure*}[!htbp]
    \centering
    \includegraphics[width=\textwidth]{figures/panel_01.png}
    \caption{\textbf{The substrate: a floor that is computed, the hypothesis it
    requires, and the cut key against ground truth.}
    \textbf{(A)}~Floor computed---not assumed---as the minimum edge weight over
    four finite contact graphs, including an adversarial instance whose smallest
    weight is $10^{-12}$. Every instance returns a strictly positive floor, because
    positivity follows from finiteness together with strict positivity of the
    weights and not from their scale.
    \textbf{(B)}~The failure branch. A system permitted to refine without limit
    generates thicknesses $1/n$ (orange): each is strictly positive and the infimum
    is zero, with every $\varepsilon$ down to $10^{-9}$ eventually cleared.
    Truncating the same sequence to a finite system (blue, first ten stages)
    restores a positive floor at $0.1$ (dashed). The uniform floor therefore
    requires finiteness, and the chart is included so that the requirement is
    visible rather than tacit.
    \textbf{(C)}~The local separation cost of \cref{eq:local-sepcost}, evaluated at
    a radius covering the graph, against an independent exhaustive enumeration of
    every bipartition, on $160$ graphs. Points lie on the identity to machine
    precision. This checks the implementation against ground truth rather than
    restating the definition.
    \textbf{(D)}~Floor surface over edge count and weight scale. The surface is
    positive everywhere on the domain and descends toward zero only as the scale
    does; increasing the edge count at fixed scale lowers the floor without
    abolishing it, which is the content of a minimum over a finite set.}
    \label{fig:panel1}
\end{figure*}


% --------------------------------------------------------------------------
\begin{figure*}[!htbp]
    \centering
    \includegraphics[width=\textwidth]{figures/panel_02.png}
    \caption{\textbf{Radius is a resolution parameter, and it runs in the direction
    opposite to the one we predicted.}
    \textbf{(A)}~Pairs of items whose derived powers agree to $10^{-9}$, against
    radius. Identification rises from zero at radius $0$ to $26$ at radius $3$. Our
    stated prediction was that a coarse radius would identify more; the measurement
    is the reverse and \cref{rem:res-inverted} records the correction.
    \textbf{(B)}~The mechanism, which is why the direction is what it is. Mean
    pairwise ball overlap (solid) rises monotonically from $0.000$ to $0.763$,
    while the number of distinct powers (dashed, right axis) falls from $138$ to
    $114$. Balls of radius $0$ are disjoint singletons and each item minimises over
    its own edges; balls of large radius coincide, so distinct items minimise over
    the same edge set and their powers converge. The overlap is plotted beside the
    count because a count without its mechanism would leave the direction
    unexplained.
    \textbf{(C)}~Six individual power trajectories on one chain as the radius
    grows. Traces that begin apart come together and then flatten, which is the
    per-item form of the aggregate in~(B).
    \textbf{(D)}~The resolution surface: derived power over item and radius. The
    surface is ridged at small radius, where items differ, and flattens along the
    radius axis as the balls merge.}
    \label{fig:panel2}
\end{figure*}


% --------------------------------------------------------------------------
\begin{figure*}[!htbp]
    \centering
    \includegraphics[width=\textwidth]{figures/panel_03.png}
    \caption{\textbf{Intensivity requires a local normaliser, and the near-miss
    control is what shows it.}
    \textbf{(A)}~Maximum drift in derived power when the graph is extended far from
    the item measured, over $120$ graphs, on a logarithmic axis. The intensive
    candidate does not move at all---the bar is at the floor of the axis because
    the value is exactly zero, not merely small. Both controls move by more than
    $10^{-1}$.
    \textbf{(B)}~The control, and the reason panel~A carries information. Fraction
    of graphs on which each candidate moved: $0/120$ for the intensive candidate
    and $120/120$ for both controls. Had no control moved, the test would have been
    measuring something other than intensivity and would be reported as
    non-discriminating.
    \textbf{(C)}~Power before against power after extension, for all three
    candidates. The intensive points lie on the identity exactly; the near-miss
    lies visibly above it, because the global floor fell and every power rose with
    it; the extensive points scatter below. The near-miss differs from the
    candidate in one component only---the normaliser---so its displacement
    localises the failure rather than merely exhibiting one.
    \textbf{(D)}~Near-miss drift over the weight of the added edge and the number
    of items added. Drift is governed by the added \emph{weight} and is nearly flat
    in the \emph{count}, which is the signature of the global floor: one small
    distant edge is enough, and further additions change little.}
    \label{fig:panel3}
\end{figure*}


% --------------------------------------------------------------------------
\begin{figure*}[!htbp]
    \centering
    \includegraphics[width=\textwidth]{figures/panel_04.png}
    \caption{\textbf{Order dependence under sequential derivation is graded, not
    abolished.}
    \textbf{(A)}~Distinct composite powers over all $720$ orderings of a six-item
    chain at radius $1$, on a logarithmic axis, with the line at $1$ marking exact
    order independence. The counts are $38$, $80$ and $432$; intensivity reduces
    dependence by an order of magnitude relative to the extensive control and no
    candidate attains $1$. The dashed line is drawn because the claim we set out to
    prove was that the intensive bar would sit on it.
    \textbf{(B)}~The same counts at radius $0$ and radius $1$. Finer locality
    tightens the intensive count from $38$ to $11$ without reaching $1$: dependence
    survives because commitment mutates edges inside the balls of neighbouring
    items, which is exactly the case the intensivity theorem does not cover.
    \textbf{(C)}~The distributions behind the counts: composite power over all
    $720$ orderings for the intensive candidate and the extensive control. The
    intensive distribution is narrow and the extensive one broad, and neither is a
    single spike.
    \textbf{(D)}~Composite power over sampled orderings and radius for the
    intensive candidate. The surface is corrugated along the ordering axis at every
    radius---that corrugation is the residual order dependence---and its amplitude
    falls as the radius falls.}
    \label{fig:panel4}
\end{figure*}


% --------------------------------------------------------------------------
\begin{figure*}[!htbp]
    \centering
    \includegraphics[width=\textwidth]{figures/panel_05.png}
    \caption{\textbf{Composition is multiplicative, and control lies at the
    strongest rung.}
    \textbf{(A)}~Mean absolute error of four candidate composition laws against a
    rung-by-rung simulation, over $4000$ ladders, on a logarithmic axis. The
    multiplicative law is exact to machine precision; the additive, maximum and
    mean laws err by $5.9\times10^{-2}$, $2.1\times10^{-1}$ and $4.6\times10^{-1}$.
    All four are scored: scoring only the favoured law would not have distinguished
    it from the others.
    \textbf{(B)}~Sensitivity normalised by the ladder's residual fraction, against
    rung power, over $700$ random ladders. Every point lies on $1/(1-\pi_j)$, which
    is exact, since sensitivity is $P/(1-\pi_j)$. The normalisation is necessary
    because raw sensitivity also depends on the other rungs, which obscures the
    relationship across ladders of differing composition.
    \textbf{(C)}~Gain in composite power from a fixed improvement applied to each
    rung in turn, at three improvement sizes, for the four-rung ladder
    $(0.45,0.30,0.55,0.20)$. Rung~3, of highest power, returns most at every size
    and rung~4, the weakest, returns least. Improving the bottleneck is the worse
    investment, by roughly a factor of two at $\delta=0.1$.
    \textbf{(D)}~Sensitivity to the third rung of a three-rung ladder, over the
    powers of the other two. The surface is $(1-\pi_1)(1-\pi_2)$: what a marginal
    gain buys is set by transmission through the remaining rungs, which is why the
    quantity being maximised is transmission and not slack.}
    \label{fig:panel5}
\end{figure*}


% --------------------------------------------------------------------------
\begin{figure*}[!htbp]
    \centering
    \includegraphics[width=\textwidth]{figures/panel_06.png}
    \caption{\textbf{The clock, the refusal, and label independence with its
    control.}
    \textbf{(A)}~The clock $M$ under the two regimes. Free rules---forming a rung,
    deriving a rung, reading a composite---leave $M$ flat at zero however many are
    performed; a climb advances it once per rung. This is the operational content
    of the freeness criterion, which licenses a free rule exactly when it commits
    no cut.
    \textbf{(B)}~The static verdict against execution on $5000$ generated
    programs: $4032$ accepted, $968$ refused, and zero disagreements, marked by the
    cross at zero. Both verdicts occur in quantity, so the test could have failed
    in either direction; a suite in which every program was accepted would not have
    tested rejection.
    \textbf{(C)}~Label independence and the control that gives it content. The
    vertical line at zero is the distribution of changes under relabelling: every
    one is exactly zero, since a rung has no identity field and the relabelling
    acts on annotations external to the formalism. The histogram is the control,
    in which a rung \emph{power} is perturbed instead; those changes are entirely
    positive and well away from zero. Without the control the first result would be
    consistent with a constant observable and would carry no information.
    \textbf{(D)}~Rungs required to reach a target, over target and rung power. The
    surface diverges as the target approaches unity and as the rung power
    approaches zero, and is low only where strong rungs meet modest targets. This
    is the cost model behind the refusal: what a target costs is set jointly by how
    demanding it is and by how much each step accomplishes.}
    \label{fig:panel6}
\end{figure*}

